\chapter{Collisions in 2D: the rehearsal}

Everything in this chapter comes back, almost unchanged, in 3D. Learn it here
where you can see it, and chapter 13 will feel like a rerun.

\section{Detecting}

Two rectangles overlap if they overlap on \textbf{both} axes at once. One axis
where they do not overlap is enough to guarantee they are apart:

\begin{center}
\begin{tikzpicture}[scale=0.82,every node/.style={font=\sffamily\footnotesize}]
  % case 1: overlap
  \begin{scope}
    \fill[rblue!25,draw=rblue] (0,0) rectangle (2.2,1.6);
    \fill[rorange!25,draw=rorange] (1.4,0.7) rectangle (3.4,2.1);
    \node at (1.7,-0.55) {overlap on X \textbf{and} Y};
    \node[text=rgreen] at (1.7,-1.05) {\textbf{HIT}};
  \end{scope}
  % case 2: no overlap on X
  \begin{scope}[xshift=5.6cm]
    \fill[rblue!25,draw=rblue] (0,0) rectangle (1.6,1.6);
    \fill[rorange!25,draw=rorange] (2.4,0.7) rectangle (4.0,2.1);
    \draw[<->,rgrey] (1.6,-0.3) -- node[below] {gap} (2.4,-0.3);
    \node at (2.0,-0.95) {no overlap on X};
    \node[text=rred] at (2.0,-1.45) {\textbf{no hit}};
  \end{scope}
\end{tikzpicture}
\end{center}

\begin{lstlisting}
bool overlapX = (a.x < b.x + b.width)  && (a.x + a.width  > b.x);
bool overlapY = (a.y < b.y + b.height) && (a.y + a.height > b.y);
bool hit = overlapX && overlapY;
\end{lstlisting}

raylib gives you that in one call, and a few relatives:

\begin{lstlisting}
CheckCollisionRecs(recA, recB);                        // box vs box
CheckCollisionCircles(centreA, rA, centreB, rB);       // circle vs circle
CheckCollisionCircleRec(centre, radius, rec);          // circle vs box
CheckCollisionPointRec(GetMousePosition(), button);    // is the mouse over it?
\end{lstlisting}

That last one is how you write a button without any UI library: a
\lstinline|Rectangle|, a \lstinline|CheckCollisionPointRec|, and
\lstinline|IsMouseButtonPressed|. Three lines, and you own it completely.

\section{Resolving, which is the hard half}

Knowing that you hit a wall does not stop you. The naive fix --- ``if I hit
something, do not move at all'' --- produces a game that feels awful: walk into a
wall at a slight angle and you stop dead instead of sliding along it.

The trick is to treat the axes separately:

\begin{center}
\begin{tikzpicture}[scale=0.9,every node/.style={font=\sffamily\footnotesize}]
  \fill[rgrey!25,draw=rgrey] (0,0) rectangle (7,0.6);
  \node[text=rgrey] at (3.5,0.3) {wall};

  \fill[rblue!30,draw=rblue] (1,0.9) rectangle (1.8,1.7);
  \draw[->,rink,line width=1pt] (1.4,1.3) -- (2.5,0.55) node[right,pos=1.1] {wanted: X and Y};

  \fill[rblue!30,draw=rblue] (3.6,0.9) rectangle (4.4,1.7);
  \draw[->,rgreen,line width=1.2pt] (4.0,1.3) -- (5.1,1.3) node[right] {X allowed};
  \draw[->,rred,line width=1.2pt,dashed] (4.0,1.3) -- (4.0,0.75) node[below right] {Y refused};
\end{tikzpicture}
\end{center}

\begin{lstlisting}
// 1. try the X move on its own
Rectangle attempt = player;
attempt.x += dx;
if (!HitsAnyWall(attempt)) player.x = attempt.x;   // otherwise: keep old x

// 2. try the Y move on its own
attempt = player;
attempt.y += dy;
if (!HitsAnyWall(attempt)) player.y = attempt.y;   // otherwise: keep old y
\end{lstlisting}

Walking diagonally into a horizontal wall, the Y step is refused and the X step
is allowed, so you slide along it. That is the whole secret, and it is the same
code in 3D with \lstinline|z| in place of \lstinline|y|.

\keyidea{Detect with one call. Resolve one axis at a time. Two moves, two
checks, never both at once.}

\section{Keeping things on screen}

The cheapest collision of all, and one every game needs:

\begin{lstlisting}
if (player.x < 0) player.x = 0;
if (player.y < 0) player.y = 0;
if (player.x + player.width  > GetScreenWidth())  player.x = GetScreenWidth()  - player.width;
if (player.y + player.height > GetScreenHeight()) player.y = GetScreenHeight() - player.height;
\end{lstlisting}

Note it clamps \emph{after} the move rather than refusing it. For a wall you want
to refuse the move; for a screen edge, clamping feels better because the player
keeps sliding along the border.

\tryit{Lesson 4 has a little maze. SPACE switches the walls between ``solid''
(resolved per axis) and ``ghost'' (detected only) so you can feel the
difference.}{./course2d/build.sh 4}
